\section{Introduction}

Neural operators have emerged as a transformative paradigm for learning solution operators of partial differential 
equations (PDEs), offering orders of magnitude speedup compared to traditional numerical solvers while maintaining 
competitive accuracy \cite{li2020fourier,lu2021deeponet}. Among neural operator architectures, Fourier Neural 
Operators (FNOs) have demonstrated remarkable success across diverse applications including fluid dynamics, 
climate modeling, and materials science. However, the deterministic nature of standard FNOs presents a critical 
limitation: they provide point predictions without quantifying prediction uncertainty. This deficiency makes 
them unsuitable for safety-critical applications such as weather forecasting, drug discovery, and engineering 
design, where reliable uncertainty estimates are essential for decision-making under uncertainty.

Uncertainty quantification (UQ) for neural networks has been extensively studied, with several prominent 
approaches emerging: Monte Carlo Dropout \cite{gal2016dropout}, Deep Ensembles \cite{lakshminarayanan2017simple}, 
and Evidential Deep Learning \cite{sensoy2018evidential,amini2020deep}. MC Dropout provides uncertainty through 
stochastic forward passes but suffers from poor calibration \cite{ovadia2019can}. Deep Ensembles offer robust 
uncertainty by aggregating independently trained models, yet incur $N$-fold computational overhead. Evidential 
methods parameterize higher-order distributions in a single forward pass, promising efficiency with explicit 
epistemic-aleatoric decomposition. However, recent studies have revealed fundamental limitations: the theoretical 
uncertainty decomposition fails in practice \cite{jurgens2024epistemic,ulmer2024prior}, and 
out-of-distribution (OOD) detection achieves near-random performance \cite{meinke2024uncertaintyquantification}. 
Despite extensive work on UQ for standard neural networks, its application to neural operators remains understudied, 
and it is unclear whether these known limitations extend to the unique context of PDE solving, characterized by 
continuous function approximation, physics-informed structure, and availability of unlimited synthetic training data.

This work presents a comprehensive empirical comparison of UQ methods for Fourier Neural Operators, 
systematically evaluating MC Dropout, Deep Ensembles, and Evidential Deep Learning across multiple dimensions: 
calibration quality, computational efficiency, uncertainty decomposition validity, and OOD detection capability. 
We implement six evidential regularization schemes and evaluate all methods on two benchmark PDEs 
(Navier-Stokes and Darcy Flow), providing the first systematic evidence base for UQ method selection in 
neural operators. The contributions of this study are summarized as follows:

\begin{enumerate}
    \item We establish a comprehensive evaluation framework for comparing UQ methods on neural operators,
    encompassing metrics across calibration (ECE, coverage, interval width), accuracy (MSE, MAE, relative $L^2$), and
    uncertainty quality (correlation, NLL, AUROC), and systematically evaluate MC Dropout, Deep Ensembles, and Evidential Deep Learning on two benchmark PDEs.

    \item We assess whether evidential uncertainty is \emph{physically} meaningful for PDE surrogates, by
    measuring the spatial alignment between predicted uncertainty and the pointwise residual of the governing
    equation---a criterion that determines whether the uncertainty is actionable for scientific and engineering
    decision-making, and which aggregate error-correlation metrics do not capture.

    \item We characterize OOD detection behavior of evidential neural operators across different types of
    distributional shift, distinguishing between parameter-level extrapolation and operator-level shift to
    understand when uncertainty-based detection is effective.

    \item We identify the conditional-independence assumption implicit in per-grid-point evidential
    parameterization as a structural limitation for operator learning, and analyze its consequences for
    domain-integrated functionals, motivating spatially correlated evidence models as the next step.
\end{enumerate}
